%MẪU CUỐN SÁCH TRONG VieTeX]
%Dùng Vietex 2.8. với phông Unicode
%Người soạn : Nguyễn Hữu Điển, ĐHKHTN, ĐHQG HN
%Mail: huudien@vnu.edu.vn, CQ: (84 - 4) 557 2869
%NR: (84 - 4) 641 8848, DĐ: 0989061951
%%%%%%%%%%%%%%%%%%%%%%%%%

\chapter{ĐỀ THI OLYMPIC IRLAND}%
\minitoc %

\thispagestyle{empty}
\Opensolutionfile{loigiaichung}[baitapC1]
\vspace*{1cm}


\section{Giới thiệu Irland}
\definecolor{xdxdff}{rgb}{0.49,0.49,1}
\definecolor{uuuuuu}{rgb}{0.27,0.27,0.27}
\definecolor{qqqqff}{rgb}{0,0,1}
\begin{tikzpicture}[line cap=round,line join=round,>=triangle 45,x=1.0cm,y=1.0cm]
\clip(-1.03,-1.01) rectangle (5.51,4.83);
\draw (4,2)-- (1,4);
\draw (1,4)-- (0,2);
\draw (0,2)-- (4,2);
\draw(2,2.25) circle (2.02cm);
\draw (1,4)-- (2.84,0.42);
\begin{scriptsize}
\fill [color=qqqqff] (0,2) circle (1.5pt);
\draw[color=qqqqff] (-0.39,2.17) node {$A$};
\fill [color=qqqqff] (4,2) circle (1.5pt);
\draw[color=qqqqff] (4.15,2.29) node {$B$};
\fill [color=qqqqff] (1,4) circle (1.5pt);
\draw[color=qqqqff] (1.15,4.29) node {$C$};
\fill [color=uuuuuu] (2,2) circle (1.5pt);
\draw[color=uuuuuu] (2.15,2.29) node {$D$};
\fill [color=xdxdff] (2.84,0.42) circle (1.5pt);
\draw[color=xdxdff] (3.01,0.69) node {$E$};
\end{scriptsize}
\end{tikzpicture}
\begin{comment}
\begin{theorem}
Tìm tất cả các cặp số nguyên $\left(x,y\right)$ sao cho $1+1996x+1998y=xy$
\end{theorem}
\end{comment}
\begin{theorem}
Tìm tất cả các cặp số nguyên $\left(x,y\right)$ sao cho $1+1996x+1998y=xy$
\end{theorem}
\definecolor{xdxdff}{rgb}{0.49,0.49,1}
\definecolor{uuuuuu}{rgb}{0.27,0.27,0.27}
\definecolor{qqqqff}{rgb}{0,0,1}
\begin{tikzpicture}[line cap=round,line join=round,>=triangle 45,x=1.0cm,y=1.0cm]
\clip(-1.25,-0.68) rectangle (5.42,4.48);
\draw [line width=1.2pt] (4,2)-- (1,4);
\draw [line width=1.2pt] (1,4)-- (0,2);
\draw [line width=1.2pt] (0,2)-- (4,2);
\draw [line width=1.2pt] (2,2.25) circle (2.02cm);
\draw [line width=1.2pt] (1,4)-- (2.84,0.42);
\begin{scriptsize}
\fill [color=qqqqff] (0,2) circle (1.5pt);
\draw[color=qqqqff] (0.13,2.26) node {$A$};
\fill [color=qqqqff] (4,2) circle (1.5pt);
\draw[color=qqqqff] (4.15,2.26) node {$B$};
\fill [color=qqqqff] (1,4) circle (1.5pt);
\draw[color=qqqqff] (1.15,4.26) node {$C$};
\fill [color=uuuuuu] (2,2) circle (1.5pt);
\draw[color=uuuuuu] (2.15,2.26) node {$D$};
\fill [color=xdxdff] (2.84,0.42) circle (1.5pt);
\draw[color=xdxdff] (2.98,0.68) node {$E$};
\end{scriptsize}
\end{tikzpicture}
% \setcounter{theorem}{7}

\glossary{name=$\mathbb{R}$,description=Kí hiệu đường thẳng thực}
\index{Hà nội}
\section{Bài tập}
\begin{baitap}
Tìm tất cả các cặp số nguyên $\left(x,y\right)$ sao cho $1+1996x+1998y=xy$

\begin{loigiai1}
Ta có: \quad $\left(x-1998\right)\left(y-1996\right)=xy-1998y-1996x+1996.1998=1997^2$\\
Do 1997 là số nguyên tố, nên ta có: $x - 1998 =  \pm 1; \pm 1997; \pm 1997^2 $. Vậy có 6 giá trị  $\left(x,y\right)$ thỏa mãn là $$ \left( {x,y} \right) = \left( {1999,1997^2  + 1996} \right),\left( {1997, - 1997^2  + 1996} \right), $$$$ \left( {3995,3993} \right),\left( {1, - 1} \right)\left( {1997^2  + 1998,1997} \right),\left( { - 1997^2  + 1998,195} \right) $$
\end{loigiai1}
\end{baitap}

\index{hanoi@Hà Nội}
\begin{baitap}
Cho $\Delta ABC$, M là điểm trong tam giác. Goi D,E,F lần lượt là hình chiếu của M xuống BC, CA, AD. Tìm tập hợp tất cả các điểm M thỏa mãn $\widehat{FDE} = \frac{\pi }{2}$
\begin{loigiai1}
Từ các tứ giác nội tiếp MDBF và MDCE ta có $\widehat{MDE} =\widehat{MCE} $ và $\widehat{MDF} =\widehat{MBE} $ do đó $\widehat{FDE} = \frac{\pi }{2} \Leftrightarrow \widehat{MCB} + \widehat{MBC} = \frac{\pi }{6}$ hay  $\widehat{BMC} = \frac{5\pi }{6}\Leftrightarrow$ M nằm trên cung tròn đi qua B và C
\end{loigiai1}
\end{baitap}
%%%%%%%%%%%%%%%%%%%%%%%%%%

\begin{baitap}
Tìm tất cả các đa thức $P(x)$ sao cho đối với mọi $x$ ta có : $$\left( {x - 16} \right)P\left( {2x} \right) = 16\left( {x - 1} \right)P\left( x \right)$$

\begin{loigiai1}

Goi $d=degP$ và $a$ là hệ số của $x$ trong $P(x)$ với số mũ lớn nhất. Khi đó hệ số của $x$ mũ lớn nhất ở bên trái là $2^d a$ phải bằng $16a$ do đó $d=4$\\
Do vế phải lúc này chia hết cho $(x-1)$, nhưng trong trường hợp đó vế phải lại chia hết cho $(x-2)$, tương tự là chia hết cho $(x-4)$ và $(x-8)$. Vậy đa thức $P(x)$ là bội của $(x-1)(x-2)(x-4)(x-8)$ là tất cả các đa thức thỏa mãn.
\end{loigiai1}
\end{baitap}
%%%%%%%%%%%%%%%%%%%%%%%%


\begin{baitap}
Cho $a,b,c $ là các số thức không âm sao cho $a + b + c \ge abc$. Chứng minh rằng $a^2  + b^2  + c^2  \ge abc$

\begin{loigiai1}
Giả sử phản chứng rằng với $a,b,c>0$ mà $a^2+b^2+c^2<abc$ do đó $abc>a^2 \Rightarrow a<bc$. Làm tương tự ta cũng có $b<ca, c<ab$. Do đó $abc\ge a^2+b^2+c^2\ge ab+bc+ca$. Theo bất đẳng thức AM-GM và $ab+bc+ca>a+b+c $ suy ra $abc>a+b+c$. Trái với giả thiết. Vậy bài toán được chứng minh.
\end{loigiai1}
\end{baitap}
%%%%%%%%%%%%%%%%%%%%%%%%%%%%%%%%%%%%

\begin{baitap}
Cho tập hợp $S = \left\{ {3,5,7,...} \right\}$. Với mỗi $x \in S$ ta đặt $\delta (x)$ là xác định một số nguyên duy nhât sao cho:\quad$2^{\delta (x)}  < x < 2^{\delta (x) + 1} $
Đối với $a,b \in S$ ta định nghĩa phép toán
$$ a\, * b = 2^{\delta (a) - 1} \left( {b - 3} \right) + a $$
a, Chứng minh rằng nếu $a,b \in S$  thì $a\, * b \in S$\\
b, Chứng minh rằng nếu $a,b,c \in S$ thì $\left( {a\, * b} \right) * c = a * \left( {b\, * c} \right)$


\begin{loigiai1}
\noindent
a, Hiển nhiên\\
b, Nếu $2^m  < a < 2^{m + 1}, 2^n  < b < 2^{n + 1}$  thì 
$$ a * b = 2^{m - 1} \left( {b - 3} \right) + a \ge 2^{m - 1} \left( {2^n  - 2} \right) + 2^m  + 1 = 2^{n + m - 1}  + 1 $$
và $a * b \le 2^{m - 1} \left( {2^{n + 1}  - 4} \right) + 2^{m + 1}  - 1 = 2^{m + n}  - 1$. Vì vậy $\delta (a * b) = m + n - 1$\\
Nếu $2^p  < c < 2^{p + 1} $ thì \\$$\left( {a * b} \right) * c = \left( {2^{m - 1} \left( {b - 3} \right) + a} \right) * c = 2^{m + n - 2} \left( {c - 3} \right) + 2^{m - 1} (b - 3) + a$$
Và$$ a * \left( {b * c} \right) = a * \left( {2^{m - 1} \left( {c - 3} \right) + b} \right) = 2^{m - 1} \left( {2^{n - 1} (c - 3) + b - 3} \right) + a = \left( {a * b} \right) * c $$

\end{loigiai1}
\end{baitap}


%%%%%%%%%%%%%%%%%%%%%%%%


\begin{baitap}
Cho tứ giác lồi ABCD có một đường tròn nội tiếp. Nếu $$A = B = \frac{{2\pi }}{3},D = \frac{\pi }{2},BC = 1$$
Tìm độ dài AD



\begin{loigiai1}
Goi I là tâm đường tròn nôi tiếp . Do $\Delta ABC$ là tam giác đều, $\widehat{BIC} = 105^0 ,\widehat{ICB} = 15^0,\widehat{AID} = 75^0,\widehat{IDA} = 45^0$ nên$$ AD = \frac{{BI}}{{BC}}\frac{{AD}}{{AI}} = \frac{{\sin 15^0 }}{{\sin 105^0 }}\frac{{\sin 75^0 }}{{\sin 45^0 }} = \sqrt 2 \sin 15^0  $$
\end{loigiai1}
\end{baitap}
%%%%%%%%%%%%%%%%%%%%%%%%


\begin{baitap}
Gọi A là tập con của $\left\{ {0,1,2,...,1997} \right\}$ gồm hơn 1000 phần tử. Chứng minh rằng A chỉ gồm những lũy thừa của 2 hoặc hai phần tử phân biệt có tổng là lũy thừa của 2



\begin{loigiai1}
\noindent
Giả  sử tập A không thỏa mãn bài toán. Khi đó A sẽ bao gồm hơn nửa số nguyên từ 51 tới 1997 mà chúng được chia thành từng cặp có tổng là 2048 $\left( {VD:51 + 1997 = 2048...} \right)$. Tương tự như vậy, A bao gồm nhiều nhất nửa số nguyên từ 14  tới 50, gồm nhiều nhất nửa số nguyên từ 3 tới 13, và có thể cả số 0, do đó A có tổng cộng  $937+18+5+1=997$ số nguyên, trái với giả thiết A gồm hơn 1000 số nguyên từ tập $\left\{ {0,1,2,...,1997} \right\}$
\end{loigiai1}
\end{baitap}
%%%%%%%%%%%%%%%%%%%%%%%%%%

\begin{baitap}
Xác định số tự nhiên n thỏa mãn những điều kiện sau:\\
a, Khai triển thập phân của n gồm 1000 số \\
b, Tất cả các số trong khai triển là số lẻ.\\
c, Hai phần tử bất kỳ liền nhau trong khai triển của n hơn kém nhau 2 đơn vị

\begin{loigiai1}
Đặt $a_n ,b_n ,c_n ,d_n ,e_n $ là số trong khai triển của n, đó là những số lẻ và hai số liên tiếp khác nhau 2 đơn vị do đó tận cùng theo thứ tự là $1,3,5,7,9$ do đó 
\[
\left( {\begin{array}{*{20}c}
   0 & 1 & 0 & 0 & 0 \\
   1 & 0 & 1 & 0 &0 \\
   0 & 1 & 0 & 1 &0  \\
   0 & 0 & 1 & 0 &1 \\
   0 & 0 & 0 & 0 &1 \\
\end{array}} \right)\left( {\begin{array}{*{20}c}
   {a_n }  \\
   {b_n }  \\
   {c_n }  \\
   {d_n }  \\
   {e_n }  \\
\end{array}} \right) = \left( {\begin{array}{*{20}c}
   {a_{n + 1} }  \\
   {b_{n + 1} }  \\
   {c_{n + 1} }  \\
   {d_{n + 1} }  \\
   {e_{n + 1} }  \\
\end{array}} \right)
\]
Gọi A là ma trận vuông trong biểu thức đó. Ta tìm giá trị riêng của của A, giả sử $Av = \lambda v$ với $v = \left( {v_1 ,v_2 ,v_3 ,v_4 ,v_5 } \right)$. Do đó 
\[
\begin{array}{l}
 v_2  = \lambda v_1  \\ 
 v_3  = \lambda v_2  - v_1  = \left( {\lambda ^2  - 1} \right)v_1  \\ 
 v_4  = \lambda v_3  - v_2  = \left( {\lambda ^3  - 2\lambda } \right)v_1  \\ 
 v_5  = \lambda v_4  - v_3  = \left( {\lambda ^4  - 3\lambda ^2  + 1} \right)v_1  \\ 
 \end{array}
\]
và $v_4  = \lambda v_5 $, do đó $\lambda ^5  - 3\lambda ^3  + \lambda  = \lambda ^3  - 2\lambda $. Giải pt này ta được $\lambda  = 0,\lambda  =  \pm 1,\lambda  =  \pm \sqrt 3 $ tương ứng ta có các vectơ riêng $x_1 ,x_2 ,x_3 ,x_4 ,x_5 $ là 
\[\left( {1,0, - 1,0,1} \right),\left( {1,1,0, - 1, - 1} \right),\left( {1, - 1,0,1, - 1} \right),\left( {1, \pm \sqrt 3 ,2, \pm \sqrt 3 ,1} \right)\]
và \[\left( {1,1,1,1,1} \right) = \frac{1}{3}x_1 \frac{{2 + \sqrt 3 }}{6}x_4  + \frac{{2 - \sqrt 3 }}{6}x_5 \]
Vì vậy \[\left( {a_{1000} ,b_{1000} ,c_{1000} ,d_{1000} ,e_{1000} } \right) = 3^{\frac{{999}}{2}} \frac{{2 + \sqrt 3 }}{6}\left( {1,\sqrt 3 ,2,\sqrt 3 ,1} \right) - \frac{{2 - \sqrt 3 }}{6}\left( {1, - \sqrt 3 ,2, - \sqrt 3 ,1} \right)\]
\[=\left( {3^{499} ,2.3^{499} ,2.3^{499} ,2.3^{499} ,3^{499} } \right)\]
Vì thế kết quả của bài toán là $8.3^{499} $
\end{loigiai1}
\end{baitap}

\Closesolutionfile{loigiaichung}
\section*{Lời giải bài tập chương 1}
\addcontentsline{toc}{section}{Lời giải bài tập chương 1}
\markright{Lời giải bài tập chương 1}
{\small\begin{Answer}{1.1}
Ta có: \quad $\left(x-1998\right)\left(y-1996\right)=xy-1998y-1996x+1996.1998=1997^2$\\
Do 1997 là số nguyên tố, nên ta có: $x - 1998 =  \pm 1; \pm 1997; \pm 1997^2 $. Vậy có 6 giá trị  $\left(x,y\right)$ thỏa mãn là $$ \left( {x,y} \right) = \left( {1999,1997^2  + 1996} \right),\left( {1997, - 1997^2  + 1996} \right), $$$$ \left( {3995,3993} \right),\left( {1, - 1} \right)\left( {1997^2  + 1998,1997} \right),\left( { - 1997^2  + 1998,195} \right) $$
\end{Answer}
\begin{Answer}{1.2}
Từ các tứ giác nội tiếp MDBF và MDCE ta có $\widehat{MDE} =\widehat{MCE} $ và $\widehat{MDF} =\widehat{MBE} $ do đó $\widehat{FDE} = \frac{\pi }{2} \Leftrightarrow \widehat{MCB} + \widehat{MBC} = \frac{\pi }{6}$ hay  $\widehat{BMC} = \frac{5\pi }{6}\Leftrightarrow$ M nằm trên cung tròn đi qua B và C
\end{Answer}
\begin{Answer}{1.3}

Goi $d=degP$ và $a$ là hệ số của $x$ trong $P(x)$ với số mũ lớn nhất. Khi đó hệ số của $x$ mũ lớn nhất ở bên trái là $2^d a$ phải bằng $16a$ do đó $d=4$\\
Do vế phải lúc này chia hết cho $(x-1)$, nhưng trong trường hợp đó vế phải lại chia hết cho $(x-2)$, tương tự là chia hết cho $(x-4)$ và $(x-8)$. Vậy đa thức $P(x)$ là bội của $(x-1)(x-2)(x-4)(x-8)$ là tất cả các đa thức thỏa mãn.
\end{Answer}
\begin{Answer}{1.4}
Giả sử phản chứng rằng với $a,b,c>0$ mà $a^2+b^2+c^2<abc$ do đó $abc>a^2 \Rightarrow a<bc$. Làm tương tự ta cũng có $b<ca, c<ab$. Do đó $abc\ge a^2+b^2+c^2\ge ab+bc+ca$. Theo bất đẳng thức AM-GM và $ab+bc+ca>a+b+c $ suy ra $abc>a+b+c$. Trái với giả thiết. Vậy bài toán được chứng minh.
\end{Answer}
\begin{Answer}{1.5}
\noindent
a, Hiển nhiên\\
b, Nếu $2^m  < a < 2^{m + 1}, 2^n  < b < 2^{n + 1}$  thì
$$ a * b = 2^{m - 1} \left( {b - 3} \right) + a \ge 2^{m - 1} \left( {2^n  - 2} \right) + 2^m  + 1 = 2^{n + m - 1}  + 1 $$
và $a * b \le 2^{m - 1} \left( {2^{n + 1}  - 4} \right) + 2^{m + 1}  - 1 = 2^{m + n}  - 1$. Vì vậy $\delta (a * b) = m + n - 1$\\
Nếu $2^p  < c < 2^{p + 1} $ thì \\$$\left( {a * b} \right) * c = \left( {2^{m - 1} \left( {b - 3} \right) + a} \right) * c = 2^{m + n - 2} \left( {c - 3} \right) + 2^{m - 1} (b - 3) + a$$
Và$$ a * \left( {b * c} \right) = a * \left( {2^{m - 1} \left( {c - 3} \right) + b} \right) = 2^{m - 1} \left( {2^{n - 1} (c - 3) + b - 3} \right) + a = \left( {a * b} \right) * c $$

\end{Answer}
\begin{Answer}{1.6}
Goi I là tâm đường tròn nôi tiếp . Do $\Delta ABC$ là tam giác đều, $\widehat{BIC} = 105^0 ,\widehat{ICB} = 15^0,\widehat{AID} = 75^0,\widehat{IDA} = 45^0$ nên$$ AD = \frac{{BI}}{{BC}}\frac{{AD}}{{AI}} = \frac{{\sin 15^0 }}{{\sin 105^0 }}\frac{{\sin 75^0 }}{{\sin 45^0 }} = \sqrt 2 \sin 15^0  $$
\end{Answer}
\begin{Answer}{1.7}
\noindent
Giả  sử tập A không thỏa mãn bài toán. Khi đó A sẽ bao gồm hơn nửa số nguyên từ 51 tới 1997 mà chúng được chia thành từng cặp có tổng là 2048 $\left( {VD:51 + 1997 = 2048...} \right)$. Tương tự như vậy, A bao gồm nhiều nhất nửa số nguyên từ 14  tới 50, gồm nhiều nhất nửa số nguyên từ 3 tới 13, và có thể cả số 0, do đó A có tổng cộng  $937+18+5+1=997$ số nguyên, trái với giả thiết A gồm hơn 1000 số nguyên từ tập $\left\{ {0,1,2,...,1997} \right\}$
\end{Answer}
\begin{Answer}{1.8}
Đặt $a_n ,b_n ,c_n ,d_n ,e_n $ là số trong khai triển của n, đó là những số lẻ và hai số liên tiếp khác nhau 2 đơn vị do đó tận cùng theo thứ tự là $1,3,5,7,9$ do đó
\[
\left( {\begin{array}{*{20}c}
   0 & 1 & 0 & 0 & 0 \\
   1 & 0 & 1 & 0 &0 \\
   0 & 1 & 0 & 1 &0  \\
   0 & 0 & 1 & 0 &1 \\
   0 & 0 & 0 & 0 &1 \\
\end{array}} \right)\left( {\begin{array}{*{20}c}
   {a_n }  \\
   {b_n }  \\
   {c_n }  \\
   {d_n }  \\
   {e_n }  \\
\end{array}} \right) = \left( {\begin{array}{*{20}c}
   {a_{n + 1} }  \\
   {b_{n + 1} }  \\
   {c_{n + 1} }  \\
   {d_{n + 1} }  \\
   {e_{n + 1} }  \\
\end{array}} \right)
\]
Gọi A là ma trận vuông trong biểu thức đó. Ta tìm giá trị riêng của của A, giả sử $Av = \lambda v$ với $v = \left( {v_1 ,v_2 ,v_3 ,v_4 ,v_5 } \right)$. Do đó
\[
\begin{array}{l}
 v_2  = \lambda v_1  \\
 v_3  = \lambda v_2  - v_1  = \left( {\lambda ^2  - 1} \right)v_1  \\
 v_4  = \lambda v_3  - v_2  = \left( {\lambda ^3  - 2\lambda } \right)v_1  \\
 v_5  = \lambda v_4  - v_3  = \left( {\lambda ^4  - 3\lambda ^2  + 1} \right)v_1  \\
 \end{array}
\]
và $v_4  = \lambda v_5 $, do đó $\lambda ^5  - 3\lambda ^3  + \lambda  = \lambda ^3  - 2\lambda $. Giải pt này ta được $\lambda  = 0,\lambda  =  \pm 1,\lambda  =  \pm \sqrt 3 $ tương ứng ta có các vectơ riêng $x_1 ,x_2 ,x_3 ,x_4 ,x_5 $ là
\[\left( {1,0, - 1,0,1} \right),\left( {1,1,0, - 1, - 1} \right),\left( {1, - 1,0,1, - 1} \right),\left( {1, \pm \sqrt 3 ,2, \pm \sqrt 3 ,1} \right)\]
và \[\left( {1,1,1,1,1} \right) = \frac{1}{3}x_1 \frac{{2 + \sqrt 3 }}{6}x_4  + \frac{{2 - \sqrt 3 }}{6}x_5 \]
Vì vậy \[\left( {a_{1000} ,b_{1000} ,c_{1000} ,d_{1000} ,e_{1000} } \right) = 3^{\frac{{999}}{2}} \frac{{2 + \sqrt 3 }}{6}\left( {1,\sqrt 3 ,2,\sqrt 3 ,1} \right) - \frac{{2 - \sqrt 3 }}{6}\left( {1, - \sqrt 3 ,2, - \sqrt 3 ,1} \right)\]
\[=\left( {3^{499} ,2.3^{499} ,2.3^{499} ,2.3^{499} ,3^{499} } \right)\]
Vì thế kết quả của bài toán là $8.3^{499} $
\end{Answer}
}